\begingroup
\shorthandoff{"<>"}
\setlength{\arrayrulewidth}{1pt}
\arrayrulecolor{vlepsbased}
\begin{longtable}{|p{15cm}|}
  \hline
  \rowcolor{maincolor} 
  \multicolumn{1}{|l|}{\textcolor{white}{\textbf{Caso de Uso}}} \\
  \hline
  \endfirsthead
  \multicolumn{1}{c}{{\tablename\ \thetable{} -- Continuación}} \\
  \hline
  \rowcolor{maincolor} 
  \multicolumn{1}{|l|}{\textcolor{white}{\textbf{Caso de Uso (Continuación)}}} \\
  \hline
  \endhead
  \hline \multicolumn{1}{|r|}{\textit{Continúa en la siguiente página...}} \\ \hline
  \endfoot
  \hline
  \caption{Caso de Uso CU6 — Cierre de sesión.}
  \label{TB:CU6_LOGOUT}
  \endlastfoot

  \rowcolor{ulepsbased}
  \begin{tabularx}{\linewidth}{@{}p{7.5cm}p{7.5cm}@{}}
    \textbf{Identificador:} CU6 & \textbf{Actor principal:} Usuario
  \end{tabularx} \\ \hline

  \textbf{Nombre:} Cierre de sesión \\ \hline

  \rowcolor{ulepsbased}
  \textbf{Descripción:} \newline 
  El usuario finaliza su actividad y cierra su sesión de forma segura revocando sus tokens de acceso. \\ \hline

  \textbf{Precondiciones:} \newline
  1. El usuario debe tener una sesión activa iniciada. \\ \hline

  \rowcolor{ulepsbased}
  \textbf{Flujo principal:} \newline
  1. El usuario autenticado pulsa "Cerrar sesión" en el menú de opciones. \newline
  2. El sistema solicita la invalidación del token de refresco en el servidor (blacklist). \newline
  3. El cliente borra los tokens del almacenamiento local. \newline
  4. El usuario es redirigido a la vista de acceso. \\ \hline

  \textbf{Flujos alternativos:} \newline
  Ninguno. \\ \hline

  \rowcolor{ulepsbased}
  \textbf{Postcondiciones:} \newline
  1. Se destruye el contexto de sesión segura. \newline
  2. El usuario pierde el acceso a las áreas privadas. \\ \hline

\end{longtable}
\endgroup
